% theme: similarity_meaning_and_ratio
\MajorQuestionLesson{相似な図形の意味と相似比}
\SetRowSpace{3mm}

\TypeQuestion{図の2つの三角形は相似である。次の問いに答えなさい。辺の長さの単位は $\mathrm{cm}$ とする。}{similarity_corresponding_triangles}


\FigProblemBlock{78mm}{fig/lesson_01_triangle_pair}{
\Text{図の2つの三角形は相似の位置にあり、対応する頂点を結ぶ線分は1点で交わっている。長さの単位は $\mathrm{cm}$ とする。}
\QQRow{頂点 $\mathrm{B}$ に対応する頂点を答えなさい。}{}{辺 $\Lseg{AC}$ に対応する辺を答えなさい。}{}
\QQRow{$\Tri{ABC}$ と $\Tri{EFD}$ の相似比を求めなさい。}{}{$\Ang{ACB}$ の大きさを求めなさい。}{}
\QQ{$\Ang{EFD}$ の大きさを求めなさい。}{}
}{
}


\TypeQuestion{図の2つの四角形は相似である。次の問いに答えなさい。辺の長さの単位は $\mathrm{cm}$ とする。}{similarity_corresponding_quadrilaterals}
\ifLMAnswerMode\else
\LMFigBlock[125mm]{fig/lesson_01_quadrilateral_pair}
\fi
\QQRow{頂点 $\mathrm{D}$ に対応する頂点を答えなさい。}{}{辺 $\Lseg{BC}$ に対応する辺を答えなさい。}{}
\QQRow{四角形 $\mathrm{ABCD}$ と四角形 $\mathrm{GHEF}$ の相似比を求めなさい。}{}{辺 $\Lseg{EF}$ の長さを求めなさい。}{}
\QQRow{$\Ang{HEF}$ の大きさを求めなさい。}{}{辺 $\Lseg{FG}$ の長さを求めなさい。}{}

\LessonQuestionBreak

\TypeQuestion{}{similarity_center_same_side}
\FigProblemBlock{78mm}{fig/lesson_01_center_same_side}{
\Text{図の2つの三角形は相似の位置にある。破線は対応する頂点を通る直線である。長さの単位は $\mathrm{cm}$ とする。}
\QQ{相似の中心を答えなさい。}{}
\QQ{辺 $\Lseg{AC}$ に対応する辺を答えなさい。}{}
\QQ{線分 $\Lseg{OD}$ の長さを求めなさい。}{}
\QQ{$\Tri{ABC}$ と $\Tri{DEF}$ の相似比を求めなさい。}{}
\QQ{$\Ang{DFE}$ の大きさを求めなさい。}{}
}{
}

\TypeQuestion{}{similarity_center_opposite_side}
\FigProblemBlock{78mm}{fig/lesson_01_center_opposite_side}{
\Text{図の2つの三角形は相似の位置にあり、対応する頂点を結ぶ線分は1点で交わっている。長さの単位は $\mathrm{cm}$ とする。}
\QQ{相似の中心を答えなさい。}{}
\QQ{頂点 $\mathrm{C}$ に対応する頂点を答えなさい。}{}
\QQ{$\Tri{ABC}$ と $\Tri{DEF}$ の相似比を求めなさい。}{}
\QQ{線分 $\Lseg{OB}$ の長さを求めなさい。}{}
\QQ{線分 $\Lseg{AD}$、$\Lseg{BE}$ の長さをそれぞれ求めなさい。}{}
}{
}
